\documentclass{beamer}
\usetheme{Boadilla}
\usecolortheme{whale}

\title[Bankruptcy Cascades]{Bankruptcy Cascades in Interbank Markets}
\subtitle{An Agent-Based Model of Systemic Risk}
\author{G. Tedeschi, A. Mazloumian, M. Gallegati, D. Helbing}
\date{2012 (Source: journal.pone.0052749)}

\begin{document}

% Slide 1: Title
\begin{frame}
    \titlepage
\end{frame}

% Slide 2: Abstract Overview
\begin{frame}{Abstract}
    \begin{itemize}
        \item Study of credit networks and interbank systems using an \textbf{Agent-Based Model (ABM)} [1].
        \item Analysis of the relationship between \textbf{business cycles} and bankruptcy cascades [1].
        \item Modeling of a three-sector economy: goods, credit, and interbank market [1].
        \item Identification of the trade-off between risk sharing and systemic risk [1].
    \end{itemize}
\end{frame}

% Slide 3: Motivation - Financialization
\begin{frame}{Motivation: Financialization of the Economy}
    \begin{itemize}
        \item In the last 30 years, the financial sector has outgrown the production sector [2].
        \item Shift from loan-based financing to speculative, market-based activities [3].
        \item Increased bank leverage and corporate debt-equity ratios have made the economy fragile [3, 4].
        \item Mainstream theory (Arrow-Debreu) fails to explain complex agent interactions during crises [5].
    \end{itemize}
\end{frame}

% Slide 4: Failure Propagation Mechanisms
\begin{frame}{Propagation of Systematic Failure}
    The literature identifies three primary types of contagion:
    \begin{enumerate}
        \item \textbf{Bank runs:} Known as self-fulfilling panics [6].
        \item \textbf{Asset price contagion:} Spirals caused by fire sales [6].
        \item \textbf{Inter-locking exposures:} Direct linkages through the interbank market [6].
    \end{enumerate}
    This paper focuses on the interbank market as a mechanism for liquidity crisis contagion [6].
\end{frame}

% Slide 5: The Risk-Sharing Trade-off
\begin{frame}{The Trade-off: Connectivity vs. Risk}
    \begin{itemize}
        \item \textbf{Benefit:} Increasing connectivity allows for \textbf{risk sharing} and diversification [7, 8].
        \item \textbf{Cost:} Excessive connectivity increases \textbf{systemic risk} by acting as a vector for contagion [8, 9].
        \item Evidence suggests a \textbf{non-monotonic relationship} between connectivity and systemic risk in the presence of agent heterogeneity [9, 10].
    \end{itemize}
\end{frame}

% Slide 6: Model Structure
\begin{frame}{Structure of the Model}
    The model consists of three interacting sectors:
    \begin{enumerate}
        \item \textbf{Goods Market:} Demand-driven output with incomplete information [11].
        \item \textbf{Credit Market:} Firms borrow from banks to finance investment [12].
        \item \textbf{Interbank Market:} Banks share risk and liquidity through random connectivity [13, 14].
    \end{enumerate}
    POPULATION: $F_t$ firms and $B_t$ banks [11].
\end{frame}

% Slide 7: Goods Market Dynamics
\begin{frame}{The Goods Market}
    \begin{itemize}
        \item Output is demand-driven; firms sell as much as the market can absorb [11].
        \item Incomplete information creates a gap between expected demand $E(D_{f,t})$ and realized demand $D_{f,t}$ [11, 15].
        \item This gap can generate unexpected shocks to revenues and negative profits [12].
    \end{itemize}
\end{frame}

% Slide 8: Firms: Production and Expectations
\begin{frame}{Firms: Production and Capital}
    \begin{itemize}
        \item Production function: $Y_{f,t} = \phi K_{f,t}$, where $\phi$ is constant capital productivity [16].
        \item Decisions follow \textbf{boundedly rational heuristics} under uncertainty [15].
        \item Firms produce based on forecasted demand; if demand exceeds capacity, they seek loans to increase capital [17, 18].
    \end{itemize}
\end{frame}

% Slide 9: Firms: Demand for Loans
\begin{frame}{Firms: Demand for Loans}
    \begin{itemize}
        \item Loan demand to reach expected demand:
        \[ L_{f,t} = \max\left\{ \frac{E(D_{f,t})}{\phi} - Y_{f,t}, 0 \right\} \quad [18] \]
        \item Adjusted for risk aversion ($\alpha$) and financial fragility:
        \[ L^d_{f,t} = \alpha \left( 1 - \frac{\Lambda L_{f,t}}{E(\pi_{f,t+1})} \right) L_{f,t} \quad [19] \]
        \item Firms decrease loan demand if they expect profits won't cover installments [19, 20].
    \end{itemize}
\end{frame}

% Slide 10: Firms: Profits and Capital Motion
\begin{frame}{Firms: Profit and Capital Motion}
    \begin{itemize}
        \item Real terms profit: $\pi_{f,t} = \min(D_{f,t}, Y_{f,t}) - \Lambda L_{f,t} \quad [20]$.
        \item All profits are retained as capital [21].
        \item Capital stock evolution:
        \[ K_{f,t} = K_{f,t-1} + \pi_{f,t} + L^d_{f,t} \quad [21] \]
    \end{itemize}
\end{frame}

% Slide 11: Banks: Interest Rate Policy
\begin{frame}{Banks: Interest Rate Policy}
    \begin{itemize}
        \item Firms contact a small number of random banks [21].
        \item Interest rate offered by bank $b$:
        \[ i^t_{f,b} = \bar{i} + c \left( \frac{L^d_{f,t}}{S_{b,t}} \right)^a \quad [22] \]
        \item $\bar{i}$ is set by the Central Bank, $S_{b,t}$ is the bank's liquidity [22].
        \item Rates decrease as the bank's \textbf{financial robustness} increases [22].
    \end{itemize}
\end{frame}

% Slide 12: Credit Rationing and Regulation
\begin{frame}{Regulatory Constraints and Rationing}
    \begin{itemize}
        \item \textbf{Basel I/II:} Regulation forces banks to hold capital caution [23].
        \item Modeled as a probability $\gamma$ of granting a loan:
        \[ \gamma^t_{f,b} = 1 - b \left( \frac{L^d_{f,t}}{S_{b,t}} \right)^s \quad [23] \]
        \item $b$ is a regulatory parameter; increasing $b$ forces higher reserves and reduces risk [23, 24].
    \end{itemize}
\end{frame}

% Slide 13: The Interbank Market Mechanism
\begin{frame}{Interbank Market Mechanism}
    \begin{itemize}
        \item If a bank lacks liquidity to cover a firm's loan, it enters the interbank market [25].
        \item It asks for funds from $x$ randomly chosen banks [25].
        \item Connectivity is modeled as a \textbf{random graph} [13].
        \item Borrower banks choose the lender offering the lowest interest rate [26].
    \end{itemize}
\end{frame}

% Slide 14: Bank Liquidity Evolution
\begin{frame}{Bank Liquidity Evolution}
    Liquidity $S_{b,t}$ evolves according to:
    \[ S_{b,t} = S_{b,t-1} - \sum L^d_{f,t} + \frac{1}{\tau} \sum L^d_{f,t'} (1 + i_{f,b}^{t'}) + I_{b,t} \quad [27] \]
    \begin{itemize}
        \item Second term: Total loans granted [27].
        \item Third term: Installments and interest from "safe" firms [27].
        \item Fourth term ($I_{b,t}$): Interbank borrowing or lending [27].
    \end{itemize}
\end{frame}

% Slide 15: Bank Profits and Bad Debt
\begin{frame}{Bank Profits}
    Profit $\pi_{b,t}$ depends on:
    \begin{enumerate}
        \item Interest paid by firms [28].
        \item \textbf{Firms' bad debt:} Losses from bankrupt firms ($\nu$ share) [28, 29].
        \item Interbank credit interactions [29].
    \end{enumerate}
    Bankruptcy occurs if the agent lacks enough cash to repay loans [30].
\end{frame}

% Slide 16: Demography of Agents
\begin{frame}{Bankruptcy and Entry Processes}
    \begin{itemize}
        \item Bankrupt agents leave the market [30].
        \item New entrants $N^{entry}_t$ are determined by expected profit opportunities, linked to the average interest rate [31, 32].
        \item Entrants are on average smaller than incumbents, mimicking empirical literature [33, 34].
    \end{itemize}
\end{frame}

% Slide 17: Simulation Results: Stylized Facts
\begin{frame}{Benchmark Model: Stylized Facts}
    The model successfully reproduces:
    \begin{itemize}
        \item \textbf{Endogenous growth} with persistent fluctuations [35].
        \item \textbf{Volatility clustering:} Large changes in variables tend to cluster together [36].
        \item Autocorrelation of absolute growth rates ($\approx 0.95$), matching G7 empirical data (0.93) [36, 37].
    \end{itemize}
\end{frame}

% Slide 18: Business Cycles and Crises
\begin{frame}{Dynamics of Business Cycles}
    \begin{itemize}
        \item Recessions are triggered by the failure of \textbf{large firms} [38].
        \item Disequilibrium dynamics lead to self-organization with persistent excess demand [38].
        \item Crises depend not on the quantity of bankruptcies, but on the \textbf{quality} (size) of bankrupted agents [39].
    \end{itemize}
\end{frame}

% Slide 19: Minsky's Financial Instability
\begin{frame}{Minsky's Financial Instability Hypothesis}
    \begin{itemize}
        \item Prosperity leads to optimism; banks grant more loans and underestimate risk [40, 41].
        \item Positive correlation (>0.63) between output and granted loans [41].
        \item Rapid increases in \textbf{financial fragility} (high leverage) anticipate severe crises [42].
    \end{itemize}
\end{frame}

% Slide 20: Size Heterogeneity
\begin{frame}{Emergent Heterogeneity}
    \begin{itemize}
        \item Trading generates \textbf{fat tail distributions} for firm and bank sizes [43].
        \item The economy is dominated by small/medium units, while a few large units control a significant part of supply [43].
        \item This makes the "representative firm" notion meaningless [44].
    \end{itemize}
\end{frame}

% Slide 21: Impact of Reserve Ratios
\begin{frame}{The Role of Reserve Requirements}
    \begin{itemize}
        \item Increasing the reserve ratio $b$ decreases bank failure rates [45].
        \item \textbf{Individual stability:} Lower average bank leverage [45].
        \item \textbf{Macroeconomic cost:} Reduces output growth by decreasing firm investment and granted loans [45, 46].
    \end{itemize}
\end{frame}

% Slide 22: Connectivity and Systemic Risk
\begin{frame}{Interbank Linkages ($x$)}
    \begin{itemize}
        \item As interbank connectivity ($x$) increases, more banks fail in each period [47].
        \item Highly connected systems exhibit \textbf{rapid survival curve declines} [48, 49].
        \item Connectivity acts as a vector for "contagious failures" [49].
        \item Higher connectivity has no effect on real output growth before cascades occur [50].
    \end{itemize}
\end{frame}

% Slide 23: Default Cascades
\begin{frame}{Large Bankruptcy Cascades}
    \begin{itemize}
        \item Measurement of the largest connected component of failed banks [51].
        \item A highly interconnected market leads to \textbf{larger cascades} of bankruptcies [51, 52].
        \item The systemic risk from connectivity outweighs the insurance benefits of risk sharing [53].
    \end{itemize}
\end{frame}

% Slide 24: Measuring "Heavy Tails"
\begin{frame}{Kurtosis and Hill Exponents}
    To measure the probability of extreme events:
    \begin{itemize}
        \item \textbf{Kurtosis:} Increases with connectivity, indicating a higher probability of systemic collapse [52].
        \item \textbf{Hill Exponent:} Decreases from "normal" values (4) toward "fragile" values (2) as interbank links increase [54].
        \item This proves that high connectivity generates \textbf{fatter tails} in failure distributions [54].
    \end{itemize}
\end{frame}

% Slide 25: Policy Implications
\begin{frame}{Conclusions and Policy}
    \begin{itemize}
        \item \textbf{Heterogeneity} is a root cause of systemic instability; failure of large banks is catastrophic [55, 56].
        \item \textbf{Policy Insight:} Interbank lending should ideally be restricted among banks with similar liquidity characteristics [57].
        \item \textbf{Redundancy:} Reserve requirements dampen cascades but limit growth [55].
    \end{itemize}
\end{frame}

% Slide 26: Summary
\begin{frame}{Final Summary}
    \begin{itemize}
        \item The model provides a tool to understand the 2008 financial meltdown [1].
        \item Confirms the \textbf{procyclicality of leverage} [58].
        \item Highlights that risk diversification across many agents has \textbf{ambiguous} effects on total systemic stability [59].
    \end{itemize}
\end{frame}

\end{document}